

\documentclass[journal, twoside]{IEEEtran}
% *** CITATION PACKAGES ***
%

%\usepackage[brazilian]{babel}%%%%% for papers in portuguese
%\usepackage[spanish]{babel}%%%%% for papers in Spanish
\usepackage{fancyhdr} %used for header
\pagestyle{fancy} %also used for header
\usepackage{lipsum}
\usepackage[utf8]{inputenc}
\usepackage[T1]{fontenc}

\usepackage{graphicx} %% for loading jpg figures
\usepackage{subfig}
\usepackage{nccmath}
\usepackage{color}
\usepackage{cite}
\usepackage{acro}
\usepackage{scalerel}
\usepackage{tikz}
\usetikzlibrary{svg.path}

\definecolor{orcidlogocol}{HTML}{A6CE39}
\tikzset{
  orcidlogo/.pic={
    \fill[orcidlogocol] svg{M256,128c0,70.7-57.3,128-128,128C57.3,256,0,198.7,0,128C0,57.3,57.3,0,128,0C198.7,0,256,57.3,256,128z};
    \fill[white] svg{M86.3,186.2H70.9V79.1h15.4v48.4V186.2z}
                 svg{M108.9,79.1h41.6c39.6,0,57,28.3,57,53.6c0,27.5-21.5,53.6-56.8,53.6h-41.8V79.1z M124.3,172.4h24.5c34.9,0,42.9-26.5,42.9-39.7c0-21.5-13.7-39.7-43.7-39.7h-23.7V172.4z}
                 svg{M88.7,56.8c0,5.5-4.5,10.1-10.1,10.1c-5.6,0-10.1-4.6-10.1-10.1c0-5.6,4.5-10.1,10.1-10.1C84.2,46.7,88.7,51.3,88.7,56.8z};
  }
}

\newcommand\orcidicon[1]{\href{https://orcid.org/#1}{\mbox{\scalerel*{
\begin{tikzpicture}[yscale=-1,transform shape]
\pic{orcidlogo};
\end{tikzpicture}
}{|}}}}



%%%% Page headears
\fancyhf{}
\fancyhead[RO]{IEEE LATIN AMERICA TRANSACTIONS ,~Vol.~7, No.~7, Oct~2020}
\fancyhead[LO]{\thepage}
\fancyhead[LE]{AUTHOR FIRST NAME  \MakeLowercase{\textit{et al.}}: MANUSCRIPT TITLE IN CAPITAL LETTERS}
\fancyhead[RE]{\thepage}

%%% header 1st page, still does not work

%\fancypagestyle{firstpage}{%
 % \lhead{IEEE LATIN AMERICA TRANSACTIONS ,~Vol.~7, No.~7, Oct~2020*}
 % \rhead{\thepage}
%}


\renewcommand{\headrulewidth}{0pt}% disable the underline of the header part

% cite.sty was written by Donald Arseneau
% V1.6 and later of IEEEtran pre-defines the format of the cite.sty package
% \cite{} output to follow that of the IEEE. Loading the cite package will
% result in citation numbers being automatically sorted and properly
% "compressed/ranged". e.g., [1], [9], [2], [7], [5], [6] without using
% cite.sty will become [1], [2], [5]--[7], [9] using cite.sty. cite.sty's
% \cite will automatically add leading space, if needed. Use cite.sty's
% noadjust option (cite.sty V3.8 and later) if you want to turn this off
% such as if a citation ever needs to be enclosed in parenthesis.
% cite.sty is already installed on most LaTeX systems. Be sure and use
% version 5.0 (2009-03-20) and later if using hyperref.sty.
% The latest version can be obtained at:
% http://www.ctan.org/pkg/cite
% The documentation is contained in the cite.sty file itself.






% *** GRAPHICS RELATED PACKAGES ***
%
\ifCLASSINFOpdf
  % \usepackage[pdftex]{graphicx}
  % declare the path(s) where your graphic files are
  % \graphicspath{{../pdf/}{../jpeg/}}
  % and their extensions so you won't have to specify these with
  % every instance of \includegraphics
  % \DeclareGraphicsExtensions{.pdf,.jpeg,.png}
\else
  % or other class option (dvipsone, dvipdf, if not using dvips). graphicx
  % will default to the driver specified in the system graphics.cfg if no
  % driver is specified.
  % \usepackage[dvips]{graphicx}
  % declare the path(s) where your graphic files are
  % \graphicspath{{../eps/}}
  % and their extensions so you won't have to specify these with
  % every instance of \includegraphics
  % \DeclareGraphicsExtensions{.eps}
\fi
% graphicx was written by David Carlisle and Sebastian Rahtz. It is
% required if you want graphics, photos, etc. graphicx.sty is already
% installed on most LaTeX systems. The latest version and documentation
% can be obtained at: 
% http://www.ctan.org/pkg/graphicx
% Another good source of documentation is "Using Imported Graphics in
% LaTeX2e" by Keith Reckdahl which can be found at:
% http://www.ctan.org/pkg/epslatex
%
% latex, and pdflatex in dvi mode, support graphics in encapsulated
% postscript (.eps) format. pdflatex in pdf mode supports graphics
% in .pdf, .jpeg, .png and .mps (metapost) formats. Users should ensure
% that all non-photo figures use a vector format (.eps, .pdf, .mps) and
% not a bitmapped formats (.jpeg, .png). The IEEE frowns on bitmapped formats
% which can result in "jaggedy"/blurry rendering of lines and letters as
% well as large increases in file sizes.
%
% You can find documentation about the pdfTeX application at:
% http://www.tug.org/applications/pdftex





% *** MATH PACKAGES ***
%
%\usepackage{amsmath}
% A popular package from the American Mathematical Society that provides
% many useful and powerful commands for dealing with mathematics.
%
% Note that the amsmath package sets \interdisplaylinepenalty to 10000
% thus preventing page breaks from occurring within multiline equations. Use:
%\interdisplaylinepenalty=2500
% after loading amsmath to restore such page breaks as IEEEtran.cls normally
% does. amsmath.sty is already installed on most LaTeX systems. The latest
% version and documentation can be obtained at:
% http://www.ctan.org/pkg/amsmath





% *** SPECIALIZED LIST PACKAGES ***
%
%\usepackage{algorithmic}
% algorithmic.sty was written by Peter Williams and Rogerio Brito.
% This package provides an algorithmic environment fo describing algorithms.
% You can use the algorithmic environment in-text or within a figure
% environment to provide for a floating algorithm. Do NOT use the algorithm
% floating environment provided by algorithm.sty (by the same authors) or
% algorithm2e.sty (by Christophe Fiorio) as the IEEE does not use dedicated
% algorithm float types and packages that provide these will not provide
% correct IEEE style captions. The latest version and documentation of
% algorithmic.sty can be obtained at:
% http://www.ctan.org/pkg/algorithms
% Also of interest may be the (relatively newer and more customizable)
% algorithmicx.sty package by Szasz Janos:
% http://www.ctan.org/pkg/algorithmicx




% *** ALIGNMENT PACKAGES ***
%
%\usepackage{array}
% Frank Mittelbach's and David Carlisle's array.sty patches and improves
% the standard LaTeX2e array and tabular environments to provide better
% appearance and additional user controls. As the default LaTeX2e table
% generation code is lacking to the point of almost being broken with
% respect to the quality of the end results, all users are strongly
% advised to use an enhanced (at the very least that provided by array.sty)
% set of table tools. array.sty is already installed on most systems. The
% latest version and documentation can be obtained at:
% http://www.ctan.org/pkg/array


% IEEEtran contains the IEEEeqnarray family of commands that can be used to
% generate multiline equations as well as matrices, tables, etc., of high
% quality.




% *** SUBFIGURE PACKAGES ***
%\ifCLASSOPTIONcompsoc
%  \usepackage[caption=false,font=normalsize,labelfont=sf,textfont=sf]{subfig}
%\else
%  \usepackage[caption=false,font=footnotesize]{subfig}
%\fi
% subfig.sty, written by Steven Douglas Cochran, is the modern replacement
% for subfigure.sty, the latter of which is no longer maintained and is
% incompatible with some LaTeX packages including fixltx2e. However,
% subfig.sty requires and automatically loads Axel Sommerfeldt's caption.sty
% which will override IEEEtran.cls' handling of captions and this will result
% in non-IEEE style figure/table captions. To prevent this problem, be sure
% and invoke subfig.sty's "caption=false" package option (available since
% subfig.sty version 1.3, 2005/06/28) as this is will preserve IEEEtran.cls
% handling of captions.
% Note that the Computer Society format requires a larger sans serif font
% than the serif footnote size font used in traditional IEEE formatting
% and thus the need to invoke different subfig.sty package options depending
% on whether compsoc mode has been enabled.
%
% The latest version and documentation of subfig.sty can be obtained at:
% http://www.ctan.org/pkg/subfig




% *** FLOAT PACKAGES ***
%
%\usepackage{fixltx2e}
% fixltx2e, the successor to the earlier fix2col.sty, was written by
% Frank Mittelbach and David Carlisle. This package corrects a few problems
% in the LaTeX2e kernel, the most notable of which is that in current
% LaTeX2e releases, the ordering of single and double column floats is not
% guaranteed to be preserved. Thus, an unpatched LaTeX2e can allow a
% single column figure to be placed prior to an earlier double column
% figure.
% Be aware that LaTeX2e kernels dated 2015 and later have fixltx2e.sty's
% corrections already built into the system in which case a warning will
% be issued if an attempt is made to load fixltx2e.sty as it is no longer
% needed.
% The latest version and documentation can be found at:
% http://www.ctan.org/pkg/fixltx2e


%\usepackage{stfloats}
% stfloats.sty was written by Sigitas Tolusis. This package gives LaTeX2e
% the ability to do double column floats at the bottom of the page as well
% as the top. (e.g., "\begin{figure*}[!b]" is not normally possible in
% LaTeX2e). It also provides a command:
%\fnbelowfloat
% to enable the placement of footnotes below bottom floats (the standard
% LaTeX2e kernel puts them above bottom floats). This is an invasive package
% which rewrites many portions of the LaTeX2e float routines. It may not work
% with other packages that modify the LaTeX2e float routines. The latest
% version and documentation can be obtained at:
% http://www.ctan.org/pkg/stfloats
% Do not use the stfloats baselinefloat ability as the IEEE does not allow
% \baselineskip to stretch. Authors submitting work to the IEEE should note
% that the IEEE rarely uses double column equations and that authors should try
% to avoid such use. Do not be tempted to use the cuted.sty or midfloat.sty
% packages (also by Sigitas Tolusis) as the IEEE does not format its papers in
% such ways.
% Do not attempt to use stfloats with fixltx2e as they are incompatible.
% Instead, use Morten Hogholm'a dblfloatfix which combines the features
% of both fixltx2e and stfloats:
%
% \usepackage{dblfloatfix}
% The latest version can be found at:
% http://www.ctan.org/pkg/dblfloatfix




%\ifCLASSOPTIONcaptionsoff
%  \usepackage[nomarkers]{endfloat}
% \let\MYoriglatexcaption\caption
% \renewcommand{\caption}[2][\relax]{\MYoriglatexcaption[#2]{#2}}
%\fi
% endfloat.sty was written by James Darrell McCauley, Jeff Goldberg and 
% Axel Sommerfeldt. This package may be useful when used in conjunction with 
% IEEEtran.cls'  captionsoff option. Some IEEE journals/societies require that
% submissions have lists of figures/tables at the end of the paper and that
% figures/tables without any captions are placed on a page by themselves at
% the end of the document. If needed, the draftcls IEEEtran class option or
% \CLASSINPUTbaselinestretch interface can be used to increase the line
% spacing as well. Be sure and use the nomarkers option of endfloat to
% prevent endfloat from "marking" where the figures would have been placed
% in the text. The two hack lines of code above are a slight modification of
% that suggested by in the endfloat docs (section 8.4.1) to ensure that
% the full captions always appear in the list of figures/tables - even if
% the user used the short optional argument of \caption[]{}.
% IEEE papers do not typically make use of \caption[]'s optional argument,
% so this should not be an issue. A similar trick can be used to disable
% captions of packages such as subfig.sty that lack options to turn off
% the subcaptions:
% For subfig.sty:
% \let\MYorigsubfloat\subfloat
% \renewcommand{\subfloat}[2][\relax]{\MYorigsubfloat[]{#2}}
% However, the above trick will not work if both optional arguments of
% the \subfloat command are used. Furthermore, there needs to be a
% description of each subfigure *somewhere* and endfloat does not add
% subfigure captions to its list of figures. Thus, the best approach is to
% avoid the use of subfigure captions (many IEEE journals avoid them anyway)
% and instead reference/explain all the subfigures within the main caption.
% The latest version of endfloat.sty and its documentation can obtained at:
% http://www.ctan.org/pkg/endfloat
%
% The IEEEtran \ifCLASSOPTIONcaptionsoff conditional can also be used
% later in the document, say, to conditionally put the References on a 
% page by themselves.




% *** PDF, URL AND HYPERLINK PACKAGES ***
%
%\usepackage{url}
% url.sty was written by Donald Arseneau. It provides better support for
% handling and breaking URLs. url.sty is already installed on most LaTeX
% systems. The latest version and documentation can be obtained at:
% http://www.ctan.org/pkg/url
% Basically, \url{my_url_here}.




% *** Do not adjust lengths that control margins, column widths, etc. ***
% *** Do not use packages that alter fonts (such as pslatex).         ***
% There should be no need to do such things with IEEEtran.cls V1.6 and later.
% (Unless specifically asked to do so by the journal or conference you plan
% to submit to, of course. )


% correct bad hyphenation here
\hyphenation{op-tical net-works semi-conduc-tor}

\DeclareAcronym{acm}{
  short = ACM ,
  long  = Association for Computing Machinery ,
  sort  = A ,
}


\usepackage{hyperref} %<--- Load after everything else
\begin{document}
 

 \setcounter{page}{7}%%%% HERE SET THE PAGE NUMBER
% \thispagestyle{firstpage}
%
% paper title
% Titles are generally capitalized except for words such as a, an, and, as,
% at, but, by, for, in, nor, of, on, or, the, to and up, which are usually
% not capitalized unless they are the first or last word of the title.
% Linebreaks \\ can be used within to get better formatting as desired.
% Do not put math or special symbols in the title.
\title{IEEE Latin America Transactions Latex Template}
%
%
% author names and IEEE memberships
% note positions of commas and nonbreaking spaces ( ~ ) LaTeX will not break
% a structure at a ~ so this keeps an author's name from being broken across
% two lines.
% use \thanks{} to gain access to the first footnote area
% a separate \thanks must be used for each paragraph as LaTeX2e's \thanks
% was not built to handle multiple paragraphs
%

\author{John Doe \orcidicon{0000-0000-0000-0000}\,, \IEEEmembership{Senior Member, IEEE}, Marco~A.~Hernandez-Nochebuena \orcidicon{0000-0000-0000-0000}\,,\IEEEmembership{ Student Member, IEEE}
        and~Ilse~Cervantes \orcidicon{0000-0000-0000-0000}\,, \IEEEmembership{Senior Member, IEEE}% <-this % stops a space
\thanks{Marco~A.~Hernandez-Nochebuena is with LINCE Lab, Instituto Politécnico Nacional Mexico e-mail:email@utfpr.edu.br.}% <-this % stops a space
\thanks{John Doe and Ilse~Cervantes are with CICATA Qro, Instituto Politécnico Nacional Mexico e-mail:email@ipn.mx}% <-this % stops a space
\thanks{Manuscript received April 19, 2025; revised August 26, 2025.}}





% The paper headers
%\markboth{IEEE Latin America Transactions,~Vol.~14, No.~8, August~2021}{Shell \MakeLowercase{\textit{et al.}}: Bare Demo of IEEEtran.cls for IEEE Journals}

%\markboth{left head}{}
%\markright{right head}
%\rhead{right head}
%\lhead{Left head}

\markboth{IEEE Latin America Transactions,~Vol.~14, No.~8, August~2021}{SKM: My IEEE article}  
% The only time the second header will appear is for the odd numbered pages
% after the title page when using the twoside option.
% 
% *** Note that you probably will NOT want to include the author's ***
% *** name in the headers of peer review papers.                   ***
% You can use \ifCLASSOPTIONpeerreview for conditional compilation here if
% you desire.




% If you want to put a publisher's ID mark on the page you can do it like
% this:
%\IEEEpubid{0000--0000/00\$00.00~\copyright~2015 IEEE}
% Remember, if you use this you must call \IEEEpubidadjcol in the second
% column for its text to clear the IEEEpubid mark.



% use for special paper notices
%\IEEEspecialpapernotice{(Invited Paper)}




% make the title area
\maketitle

% As a general rule, do not put math, special symbols or citations
% in the abstract or keywords.
\begin{abstract}

\lipsum[2-3]
\end{abstract}

% Note that keywords are not normally used for prerreview papers.
\begin{IEEEkeywords}
Latex template, IEEE, Latin America Transactions, guidelines for authors.
\end{IEEEkeywords}







\IEEEpeerreviewmaketitle



\section{Introdução}

\IEEEPARstart{E}{m} diversas aplicações \lipsum[1-4] \cite{bodaghi2017passive,ribeiro2015passive}, \cite{koroishi2015robust,koroishi2016robust}. \cite{min2017concept,mohammadi2015modeling},
% \cite{koroishi2015robust,koroishi2016robust,yang2018networked},  \cite{liu2018dynamic,liang2017optimal},  \cite{luo2018active,koroishi2014numerical}  \cite{nugraha2018setting} 


 

\lipsum[1-5]
\ac{acm}, is an abbreviation. \Ac{acm} is an association.


\section{Metodos}
\lipsum[1-5]
 Fig.~\ref{fig:est}. A plataforma experimental apresenta-se na Fig.~\ref{fig:bancada} e o diagrama para obtenção da equação de movimento apresenta-se na Fig.~\ref{fig:diagrama}.


\begin{figure}[!t]
\centering
\subfloat[Bancada]{\includegraphics[width=2.5in]{figure/img2.png}%
\label{fig:bancada}}
\hfil
\subfloat[Diagrama esquematico]{\includegraphics[width=2.5in]{figure/img4.png}%
\label{fig:diagrama}}
\caption{Estrutura de dois graus de liberdade.}
\label{fig:est}
\end{figure}

A equação de movimento da estrutura apresenta-se na Eq. (\ref{eq:1}).
\begin{equation}
	[M]\{\ddot{x} \}+[C]\{\dot{x} \}+[K]\{x \} = \{F(t) \} +\{F_{EAM}(t) \}
	\label{eq:1}
\end{equation}
Sendo $[M]$, $[K]$ e $[C]$ as matrizes de massa, rigidez e amortecimento, respectivamente. $\{F\}=\{F_1(t) \quad F_2(t)\}^T$ é o vetor da força de excitação, $\{F_{AEM}\}$ o vetor de força do atuador eletromagnético na massa $m_1$ e $\{x(t)\}=\{x_1(t) \quad x_2(t)\}^T$ o vetor deslocamento.

\begin{figure*}[h] 
   \caption{Animalitos bonitos}\label{figur:daily_instantaneous}
  %\begin{center}
 \subfloat[SOC(t) CNCQR  ]{\label{figur:SOCCNCQR}\includegraphics[width=.33\textwidth]{figure/img1.png}} %%
 \subfloat[SOC(t) DIABN ]{\label{figur:socdia}\includegraphics[width=.33\textwidth]{figure/img3.png}}
 \subfloat[SOC(t) TEZDF  ]{\label{figur:soctez}\includegraphics[width=.33\textwidth]{figure/img2.png}}
  \\
 \subfloat[H$_2$ mass CNCQR  ]{\label{figur:h2CNC}\includegraphics[width=.33\textwidth]{figure/img4.png}}
 \subfloat[H$_2$ mass DIABN  ]{\label{figur:h2DIA}\includegraphics[width=.33\textwidth]{figure/img5.png}}
 \subfloat[H$_2$ mass TEZDF  ]{\label{figur:h2tez}\includegraphics[width=.33\textwidth]{figure/img6.png}}
 %  \end{center}
   \end{figure*}

      
      
\begin{table*}[h]
\caption{Daily average load and  minimum design requirements as a function of the location and the season of the year.}
\label{table_minimumsurface}
\centering
\resizebox{18cm}{!} {
\begin{tabular}{cccccccccc}
\hline
\hline
	&		&		&		\multicolumn{2}{c}{$\int_{24h} P dt$\hspace{2mm} [units]}	& Column	&	Column &	Columns	&	Column	&	Column 		\\
Location	&	Season	&	Load	&[Units]	&		& [Units]	&  [Units]	&	 [Units]&	[Units] 	&	[Units] \\
\hline
Name1	&	type A	&	Spring	&	9.40	&	0.28	&	3.70	&	16.8	&	6	&	35	&	0	\\
	&	type A	&	Summer	&	9.20	&	0.38	&	3.50	&	16.8	&	6	&	35	&	0	\\
	&	type A	&	Fall	&	8.56	&	0.53	&	2.86	&	16.8	&	6	&	35	&	3	\\
	&	type A	&	Winter	&	8.50	&	0.63	&	2.79	&	16.8	&	6	&	35	&	4	\\
	&	type B	&	Spring	&	20.70	&	0.25	&	3.55	&	57.6	&	23	&	120	&	0	\\
	&	type B	&	Summer	&	20.45	&	0.41	&	3.29	&	57.6	&	23	&	120	&	0	\\
	&	type B	&	Fall	&	20.02	&	0.64	&	2.87	&	57.6	&	23	&	120	&	6	\\
	&	type B	&	Winter	&	19.98	&	0.64	&	2.83	&	57.6	&	23	&	120	&	5	\\
	Name2	&	type A	&	Spring	&	9.40	&	0.28	&	3.70	&	16.8	&	6	&	35	&	0	\\
	&	type A	&	Summer	&	9.20	&	0.38	&	3.50	&	16.8	&	6	&	35	&	0	\\
	&	type A	&	Fall	&	8.56	&	0.53	&	2.86	&	16.8	&	6	&	35	&	3	\\
	&	type A	&	Winter	&	8.50	&	0.63	&	2.79	&	16.8	&	6	&	35	&	4	\\
	&	type B	&	Spring	&	20.70	&	0.25	&	3.55	&	57.6	&	23	&	120	&	0	\\
	&	type B	&	Summer	&	20.45	&	0.41	&	3.29	&	57.6	&	23	&	120	&	0	\\
	&	type B	&	Fall	&	20.02	&	0.64	&	2.87	&	57.6	&	23	&	120	&	6	\\
	&	type B	&	Winter	&	19.98	&	0.64	&	2.83	&	57.6	&	23	&	120	&	5	\\
\hline
\hline
\end{tabular}}
\end{table*}
      
      
      
      

\section{Results}

\lipsum[5-10]  \cite{koroishi2016robust}, \cite{ahlawat2001multiobjective}.
 A função objetivo é apresentada pela Eq.~(\ref{eq:2}), sendo $FRF_{exp}$ a $FRF$ experimental e $FRF_{simu}$ a simulada.

\begin{equation}
	FO=\sum \left| \frac{FRF_{exp}- FRF_{simu}}{FRF_{exp}} \right|
	\label{eq:2}
\end{equation}

\begin{table}[!t]
	\caption{Los parametros.}
	\label{tab:FO}
	\centering
		\begin{tabular}{lllll} %%%% DO NOT USE VERTICAL LINES
		\hline
	\hline
	&	Columna 1	&	Columna 2	&	Columna 3	&	Columna 4	\\
	\hline
Renglón 1	&	3.33	&	38.4	&	15	&	80	\\
Renglón 2	&	3.56	&	16.8	&	6	&	35	\\
Renglón 3	&	3.07	&	16.8	&	6	&	35	\\
Renglón 4	&	2.94	&	16.8	&	6	&	35	\\
Renglón 5	&	2.94	&	16.8	&	6	&	35	\\
Renglón 6	&	3.95	&	40.8	&	16	&	85	\\
Renglón 7	&	3.34	&	40.8	&	16	&	85	\\
Renglón 8	&	3.17	&	40.8	&	16	&	85	\\
Renglón 9	&	3.27	&	40.8	&	16	&	85	\\
		\hline
	\hline
	\end{tabular}%}
\end{table}

A Fig.~\ref{figur:daily_instantaneous} apresenta a comparação entre a FRF experimental e as FRFs identificadas com o menor valor da função objetivo para os diferentes valores de população considerando as quatro FRFs segundo as entradas e saídas definidas na equação de movimento do modelo apresentado na Eq. (\ref{eq:1}): $x_2/F_1$ (ver a Fig.~\ref{figur:h2DIA}), $x_1/F_1$ (ver a Fig.~\ref{figur:h2CNC}), $x_1/F_2$ (ver a Fig.~\ref{figur:h2tez}), $x_2/F_2$ (ver a Fig.~\ref{figur:SOCCNCQR}). \cite{ahlawat2001multiobjective}

\lipsum[11-15]

\section{Conclusiones}

\lipsum[5-10]









      
      






% if have a single appendix:
%\appendix[Proof of the Zonklar Equations]
% or
%\appendix  % for no appendix heading
% do not use \section anymore after \appendix, only \section*
% is possibly needed

% use appendices with more than one appendix
% then use \section to start each appendix
% you must declare a \section before using any
% \subsection or using \label (\appendices by itself
% starts a section numbered zero.)
%

% use section* for acknowledgment
\section*{Agradecimentos}

Os autores agradecem o Conselho Nacional de Desenvolvimento Científico e Tecnológico (CNPq) (Processo xxxxx/20xx-4), a Fundação Araucária e a Coordenação de Aperfeiçoamento de Pessoal de Nível Superior (CAPES) pelos suportes financeiros para o desenvolvimento do presente trabalho.




\bibliographystyle{ieeetr}  %%%% https://www.overleaf.com/learn/latex/bibtex_bibliography_styles

\bibliography{References}

% Can use something like this to put references on a page
% by themselves when using endfloat and the captionsoff option.
\ifCLASSOPTIONcaptionsoff
  \newpage
\fi


%\begin{thebibliography}{1}

%\bibitem{IEEEhowto:kopka}
%H.~Kopka and P.~W. Daly, \emph{A Guide to \LaTeX}, 3rd~ed.\hskip 1em plus
%  0.5em minus 0.4em\relax Harlow, England: Addison-Wesley, 1999.
%
%\end{thebibliography}






\begin{IEEEbiography}[{\includegraphics[width=1in,height=1.25in,clip,keepaspectratio]{figure/unknowngirl.png}}]{2nd author}
\lipsum[1-2]
\end{IEEEbiography}

% if you will not have a photo at all:
\begin{IEEEbiography}[{
\includegraphics[width=1in,height=1.25in,clip,keepaspectratio]{figure/unknowngirl.png}
}]
{2nd author}
\lipsum[1]
\end{IEEEbiography}

% insert where needed to balance the two columns on the last page with
% biographies
%\newpage

\begin{IEEEbiography}[{\includegraphics[width=1in,height=1.25in,clip,keepaspectratio]{figure/unknownman.png}}]{3rd author}
\lipsum[1]
\end{IEEEbiography}
\vspace{11cm}
\printacronyms 
\end{document}


